\chapter{Methodology}

This chapter describes the experimental methodology %szb: sztem nem kell experimental csak methodology
used to evaluate deep learning-based code translation models. The study focuses
on comparing models under a unified experimental setup, using both
similarity-based and execution-based evaluation techniques.

Code translation outputs are evaluated using BLEU score as a similarity-based
metric, while functional correctness %szb: mintha szembeállítanád, mintha ez nem vonatkozna a Code translation outputs kategóriába. De igen, hisz ezekkel is ezt értékeled
is assessed with the CodeEditorBench framework, which assesses model performance
on practical coding tasks such as translation and debugging, providing a more
realistic measure of functional correctness.

The following sections detail the experimental setup, models, dataset, translation %szb: jelzők jöhetnének: milyen models (finethined), milyen dataset? (training) ilyesmi
pipeline, and evaluation procedure.

\section{Experimental Setup}

The experiments were designed to compare the performance of base and fine-tuned
code translation models under identical conditions. Two open-source
code-specific large language models were selected for evaluation: the
\texttt{bigcode/starcoder2-3b} model \cite{starcoder2} and the
\texttt{Qwen/Qwen2.5-Coder-3B-Instruct} model \cite{qwen_coder}.

For both models, two experimental variants were prepared. The first
variant used the original pre-trained model %szb: ez így megtévesztő, ugyanis az általad használt Qwen az már finetuneolt modell, így kicsit fura original pre-trained-ként jellemezni. Ink azt írd, hogy original models without any custom/specific funetuning/posttraining for code-transformation, vagy vmi ilyesmi, és akk egyértelmű, hogy itt a te saját finetuningodra gondolsz
without additional task-specific adaptation. The second variant used a
fine-tuned version of the same model, trained specifically for Java-to-C++ code
translation. This setup enables a direct comparison between base %szb: ugyanígy, ne hívd base modellnek, mert az sokszor azt jelenti, hogy a nem instrukciókövetésr ehangolt modell
and fine-tuned models, making it possible to analyze the impact of task-specific
adaptation on translation quality.


The experiments focus exclusively on Java-to-C++ translation. This direction was
selected because Java and C++ represent two widely used but substantially
different programming languages. Java relies on automatic garbage collection,
the Java Virtual Machine and higher-level abstractions, while C++ requires
manual memory management and provides lower-level control over system resources.
These differences make Java-to-C++ translation more challenging than translation
between more closely related languages. In addition, the XLCoST dataset used for
fine-tuning in this thesis contains aligned Java and C++ code pairs, and the
CodeEditorBench framework supports evaluation of this language pair, allowing
consistent comparison across all models.

The experimental workflow consisted of two main phases. In the first phase, the
base models were evaluated directly on the Java-to-C++ translation task in order
to measure their original performance without additional task-specific
adaptation. For each Java input sample, the model generated a C++ translation.
The generated output was then compared against the reference
\texttt{target\_code} provided in the CodeEditorBench dataset in order to
compute BLEU score. The same generated translation was also compiled and
executed inside the CodeEditorBench environment to assess functional
correctness.

In the second phase, the same models were fine-tuned for Java-to-C++ code
translation and evaluated again using the identical procedure. This made it
possible to directly compare the performance of the base and fine-tuned
variants.

The XLCoST dataset \cite{zhu2022xlcost} %szb: elég csak első említéskor hivatkozni - ez mindenre is igaz
was used as the primary training resource for the fine-tuned models. XLCoST
contains aligned code snippets across multiple programming languages, including
Java and C++, making it suitable for supervised code translation tasks. The
Java--C++ pairs were first extracted and transformed into instruction-based
prompts, then tokenized and prepared for training. Fine-tuning was performed
using parameter-efficient methods in order to reduce hardware requirements and
training time while still allowing the models to adapt to the translation task.

%szb: ez a kövi 2 bekezdés mehetne 1be, és akkor valahogy úgy bevezetni inkább, hogy "As mentioned above, one of the two models selected for the experiments is StarCoder, ... which was selected"
The StarCoder2-3B model \cite{starcoder2, starcoder2_hf} was selected because it
provides strong coding capabilities while remaining relatively lightweight
compared to larger code models. StarCoder2 models are trained on trillions of
tokens collected from large-scale code corpora. The 3B variant was trained on
more than 3 trillion tokens and, at the time of its release, outperformed
several other code models of similar size.

The Qwen2.5-Coder-3B-Instruct model was selected because it is specifically
optimized for programming-related tasks and has demonstrated strong results on
code generation and repair benchmarks \cite{qwen_coder}. The Qwen2.5-Coder
series is trained on more than 5.5 trillion tokens and achieves competitive
performance across a wide range of coding tasks \cite{qwen_coder,
qwen_coder_report}.

Execution-based evaluation was performed using the CodeEditorBench framework.
CodeEditorBench evaluates model outputs in realistic coding scenarios,
including translation tasks, by compiling and running generated code against
predefined test cases \cite{codeeditorbench}. The evaluation environment was
executed inside a Docker container using the official CodeEditorBench image.
Using Docker ensured that all experiments were performed in an isolated and
reproducible environment, reducing the risk of issues caused by missing
dependencies, compiler version differences, or operating system variations.
Container-based execution also provides sandboxing, which is important when
running automatically generated code because it limits the potential impact of
compilation failures, runtime errors or unintended system-level behavior.


To ensure fair comparison, all models were evaluated on the same Java-to-C++
pairs of CodeEditorBench and used the same instruction format,
preprocessing pipeline, and generation settings. Each input example was
converted into an instruction-following prompt containing a system message, a
user request, the original Java code, and a response tag indicating where the
C++ translation should begin.

The prompt format used during both training and inference was the following:



\lstset{
  basicstyle=\ttfamily\small,
  breaklines=true,
  breakatwhitespace=true,
  columns=fullflexible
}

\begin{lstlisting}
### System
You are an expert code translator. Translate Java code to C++ preserving logic and I/O.

### User
Translate this Java code to C++.

### Java
<source code>

Answer (C++)
\end{lstlisting}

Only the generated C++ part after the \texttt{\#\#\# Answer (C++)} tag was used
as the training target. During preprocessing, all prompt tokens before the
response tag were masked with \texttt{-100} labels so that the loss was
calculated only on the translated C++ output.

The same tokenizer settings, truncation strategy, and maximum sequence length
of 1024 tokens were used for all models. During inference, generation was
performed with a maximum of 900 new tokens, a repetition penalty of 1.10,
\texttt{top\_p=0.95}, and \texttt{temperature=0.2}. Greedy decoding was used
with \texttt{do\_sample=False} in order to ensure deterministic outputs across
repeated runs. %szb: picit fura, ha nemnulla a temperature, akkor a do\_sample true szokott lenni. De mostmár ne értékeld persze újra ki az egészet... De lehet, hogy kivenném a do_sample-t, nehogy kérdéseket vessen fel.

After generation, all outputs were postprocessed using the same cleaning
pipeline. This step removed markdown code fences, duplicated prompt fragments,
incomplete trailing sections and any repeated system or user messages. BLEU
score was then computed against the reference C++ translation, while
execution-based evaluation was performed by compiling and running the generated
code inside the Docker-based CodeEditorBench environment.